\subsection{Semidirect Products}

Let $H$ and $K$ be groups, let $\phi : K \to \aut(H)$ be a homomorphism, and identify $H, K \leq G = H \semidir_\phi K$.

\begin{exercise}[5.5.1]
    Prove that $C_K(H) = \ker\phi$ (recall that $C_K(H) = C_G(H) \cap K$).
\end{exercise}

\begin{sol}
    Let $k \in C_K(H)$. Then $k \in K$ and $kh = hk$ for all $h \in H$. Then $\phi(k)(h) = khk\inv = h$ for all $h \in H$, so $k \in \ker\phi$. Conversely, let $k \in \ker\phi$. Then $\phi(k)(h) = khk\inv = h$ for all $h \in H$, so $kh = hk$ for all $h \in H$, hence $k \in C_K(H)$. Therefore, $C_K(H) = \ker\phi$.
\end{sol}

\begin{exercise}[5.5.2]
    Prove that $C_H(K) = N_H(K)$.
\end{exercise}

\begin{sol}
    It is clear that $C_H(K) \leq N_H(K)$. Now let $h \in N_H(K)$. Then $hKh\inv = K$, so for any $k \in K$, we have $hkh\inv \in K$. Since $H \cap K = 1$, and $[h, k] = hkh\inv k\inv \in H \cap K$, we must have $[h, k] = 1$, so $hk = kh$ for all $k \in K$. Hence, $h \in C_H(K)$, so $N_H(K) \leq C_H(K)$. Therefore, $C_H(K) = N_H(K)$.
\end{sol}

\begin{exercise}[5.5.3]
    In Example 1 following the proof of Proposition 5.11, prove that every element of $G - H$ has order 2. Prove that $G$ is Abelian if and only if $h^2 = 1$ for all $h \in H$.
\end{exercise}

\begin{sol}
    Recall that $H \nsub G$. It is clear to see that $G/H = \set{H, Hk}$ for $k \in H$, since $H$ consists of elements $hk$ with trivial $k$ component. Then for any $hk \in Hk$, and noting the fact that $kh = h\inv k$ for all $h \in H$, we have
    \[(hk)^2 = h(kh)k = h(h\inv k)k = kk = 1.\]
    If $hk = 1$, then $h = k\inv \in H \cap K = 1$, so $hk = 1$ if and only if $h = k = 1$. Hence, every element of $G - H$ has order 2.

    If $G$ is Abelian, then for any $h \in H$ and $k \in K$, we have $hk = kh$. Moreover, $kh = h\inv k$, hence $hk = h\inv k$, so $h^2 = 1$. Conversely, if $h^2 = 1$ for all $h \in H$, then any $h \in H$ and $k \in K$ commute, because $kh = h\inv k = hk$ by $h^2 = 1$. Then for any $h_1, h_2 \in H$ and $k_1, k_2 \in K$, we have
    \[(h_1k_1)(h_2k_2) = h_1(h_2k_1)k_2 = (h_2h_1)(k_2k_1) = (h_2k_2)(h_1k_1),\]
    so $G$ is Abelian.
\end{sol}

\begin{exercise}[5.5.4]
    Let $p = 2$ and check that the construction of the two non-Abelian groups of order $p^3$ is valid in this case. Prove that \emph{both} resulting groups are isomorphic to $D_8$.
\end{exercise}

\begin{sol}
    Let $U = \gen u \cong Z_4$ and $V = \gen v \cong Z_2$. Then $\aut(U) \cong Z_2$, so there is a unique non-trivial homomorphism $\phi : V \to \aut(U)$, which sends $v$ to the automorphism of $U$ that sends $u$ to $u\inv$. Then the semidirect product $U \semidir_\phi V$ has order 8, and is generated by $u$ and $v$ with relations
    \[u^4 = v^2 = 1, \quad vu = u\inv v.\]
    It is clear that $U \semidir_\phi V$ is non-Abelian, since $vu \neq uv$. Let $r = u$ and $s = v$. Then $r^4 = s^2 = 1$, and $sr = vu = u\inv v = r\inv s$, so $U \semidir_\phi V$ has the same presentation as $D_8$, hence $U \semidir_\phi V \cong D_8$.

    Now let $W = \gen w \times \gen x \cong V_4$ and $Y = \gen y \cong Z_2$. Then $\aut(W) \cong S_3$. Since $S_3$ contains 3 transpositions that are conjugate to each other, choose some $\phi : Y \to \aut(W)$ that sends $y$ to a transposition in $\aut(W)$. Then the semidirect product $W \semidir_\phi Y$ has order 8, and is generated by $w$, $x$, and $y$ with relations
    \[w^2 = x^2 = y^2 = 1, \quad yw = x y, \quad yx = w y.\]
    It is clear that $W \semidir_\phi Y$ is non-Abelian, since $yw \neq wy$. Let $r = yw$ and $s = y$. Then $s^2 = 1$,
    \[r^4 = (r^2)^2 = ((yw)^2)^2 = (xyyw)^2 = (xw)^2 = x^2w^2 = 1,\]
    and
    \[sr = y(yw) = w = w\inv = w\inv y\inv y = (yw)\inv y = r\inv s,\]
    so $W \semidir_\phi Y$ has the same presentation as $D_8$, hence $W \semidir_\phi Y \cong D_8$.
\end{sol}

\begin{exercise}[5.5.5]
    Let $G = \hol(Z_2 \times Z_2)$.
    \begin{subexercises}
        \item Prove that $G = H \semidir K$, where $H = Z_2 \times Z_2$ and $K \cong S_3$. Deduce that $|G| = 24$.
        \item Prove that $G$ is isomorphic to $S_4$. [Obtain a homomorphism from $G$ into $S_4$ by letting $G$ act on the left cosets of $K$. Use \exref{5.5.1} to show that this representation is faithful.]
    \end{subexercises}
\end{exercise}

\begin{solenum}
    \item Note that $\aut(H) \cong \gl_2(\f_2) \cong S_3$ as it is a non-Abelian group of order 6, by \exref{4.2.10}. It follows by definition that $G = H \semidir K$, where $K \cong S_3$. Hence, $|G| = |H||K| = 4 \cdot 6 = 24$.
    \item Let $G$ act on the left cosets of $K$ in $G$. From $|G : K| = 4$, we have a homomorphism $\phi : G \to S_4$. Moreover, we have some homomorphism $\psi : K \to \aut(H)$ such that $G \cong H \semidir_\psi K$. Lastly, we know from Theorem 4.3 that
    \[\ker\phi = \bigcap_{g \in G} gKg^{-1} \nsub G\]
    such that $\ker\phi \leq K$. Putting $N = \ker\phi$, then $N = \set{1, A_3, K}$. If $N = K$, then $K \nsub G$, which is a contradiction since $K$ is not normal in $G$ because $\psi$ is injective. If $N = A_3$, let $n \in N$ be a non-identity element. Then $hnh\inv \in N$ for every $h$. However, \exref{5.5.1} says that $\ker\psi = C_K(H)$. Since $\psi$ is injective (an isomorphism, actually), then $C_K(H) = 1$, so $n$ does not commute with some $h_0 \in H$, i.e., $nh_0n\inv \neq h_0$. However, from $N \nsub G$, we know that $h_0nh_0\inv \in N$. 
    
    Observe that $h_0nh_0\inv = (h_0nh_0\inv n\inv)n = [h_0, n]n$. Observe that $h_0(nh_0\inv n\inv) \in H$ by definition of the semidirect product, so $[h_0, n] \in H$. If $[h_0, n] = 1$, then $h_0nh_0\inv = n \in N$, contradicting our choice of $h_0$. Hence $h_0nh_0\inv$ has non-trivial component in $H$, so $h_0nh_0\inv \notin N$, which is a contradiction. Therefore, $N = 1$, so $\phi$ is injective, hence $G \cong \phi(G) \leq S_4$. Since $|G| = |S_4| = 24$, we have $G \cong S_4$.
\end{solenum}

\begin{spex}[5.5.6]
    Assume that $K$ is a cyclic group, $H$ is an arbitrary group, and $\phi_1$ and $\phi_2$ are homomorphisms from $K$ into $\aut(H)$ such that $\phi_1(K)$ and $\phi_2(K)$ are conjugate subgroups of $\aut(H)$. If $K$ is infinite, assume $\phi_1$ and $\phi_2$ are injective. Prove, by constructing an explicit isomorphism, that
    \[H \semidir_{\phi_1} K \cong H \semidir_{\phi_2} K\]
    (in particular, if the subgroups $\phi_1(K)$ and $\phi_2(K)$ are equal in $\aut(H)$, then the resulting semidirect products are isomorphic). [Suppose $\sigma\phi_1(K)\sigma\inv = \phi_2(K)$ so that for some $a \in \z$ we have $\sigma\phi_1(k)\sigma\inv = \phi_2(k)^a$ for all $k \in K$. Show that the map
    \[\psi : H \semidir_{\phi_1} K \to H \semidir_{\phi_2} K\]
    defined by $\psi((h,k)) = (\sigma(h), k^a)$ is a homomorphism. Show $\psi$ is bijective by constructing a 2-sided inverse.]
\end{spex}

\begin{sol}
    Let $K = \gen k$ and $\sigma \in \aut(H)$ such that $\sigma\phi_1(K)\sigma\inv = \phi_2(K)$. Then there exists some $a \in \z$ such that $\sigma\phi_1(k)\sigma\inv = \phi_2(k)^a$ for all $k \in K$. Define the map
    \[\psi : H \semidir_{\phi_1} K \to H \semidir_{\phi_2} K\]
    by $\psi((h, k^m)) = (\sigma(h), k^{am})$ for all $h \in H$ and $m \in \z$. Then for any $(h_1, k^s), (h_2, k^t) \in H \semidir_{\phi_1} K$, we have
    \begin{align*}
        \psi((h_1, k^s)(h_2, k^t)) & = \psi((h_1\phi_1(k^s)(h_2), k^{s + t})) \\
        & = (\sigma(h_1\phi_1(k^s)(h_2)), k^{a(s + t)}) \\
        & = (\sigma(h_1)\sigma(\phi_1(k^s)(h_2)), k^{as}k^{at}) \\
        & = (\sigma(h_1)\phi_2(k^{as})(\sigma(h_2)), k^{as}k^{at}) \\
        & = (\sigma(h_1), k^{as})(\sigma(h_2), k^{at}) \\
        & = \psi((h_1, k^s))\psi((h_2, k^t)),
    \end{align*}
    so $\psi$ is a homomorphism. To see that $\psi$ is bijective, we split into two cases.
    \begin{enumerate}
        \item Suppose $K$ is infinite. By conjugate subgroups, we know that $\phi_2(K) = \phi_2(K^a)$, where $K^a = \set{k^a \mid k \in K}$. Since $\phi_2$ is injective, then $K^a = K$, forcing $a = \pm 1$. Then $\theta((h, k)) = (\sigma\inv(h), k^a)$ is a 2-sided inverse of $\psi$, so $\psi$ is bijective.
        \item We first prove a lemma: if $m \mid n$, then the map defined by
        \[\nu : \units[n] \to \units[m], \quad \nu(a) = a \bmod m\]
        is a surjective homomorphism. We first check that $\nu$ is well-defined: if $a \equiv b \bmod n$, then $n \mid (a - b)$, so $m \mid (a - b)$, hence $a \equiv b \bmod m$. Next, we check that $\nu$ is a homomorphism: for any $a, b \in \units[n]$, we have
        \[\nu(ab) = ab \bmod m = (a \bmod m)(b \bmod m) = \nu(a)\nu(b).\]
        Finally, we check that $\nu$ is surjective: let $a \in \units[m]$. If $a \in \units[n]$, then $\nu(a) = a$. If $a \notin \units[n]$, then $(a, n) \neq 1$. Let $c$ be the product of primes that divide $n$ but not $a$, let $p$ be a prime that divides $n$, and put $b = a + cm$. We claim that $b \in \units[n]$ as follows:
        \begin{itemize}
            \item If $p \mid a$, then $p \nmid c$ by definition of $c$, hence $p \nmid b$.
            \item If $p \nmid a$, then $p \mid c$ by definition of $c$, hence $p \nmid b$.
        \end{itemize}
        Thus, $b \in \units[n]$, and $\nu(b) = b \bmod m = a + cm \bmod m = a$. Hence, $\nu$ is surjective.
        
        Suppose $K \cong Z_n$ is finite. Let $|\phi_1(K)| = |\phi_2(K)| = m$, where $m \mid n$. Note that $\gen{\phi_1(k)} = \phi_1(K)$ and $\gen{\phi_2(k)} = \phi_2(K)$, so $|\phi_1(k)| = |\phi_2(k)| = m$. Moreover, $\gen{\phi_2(k)^a} = \phi_2(K)$, so $|\phi_2(k)^a| = m$. It follows that $(a, m) = 1$. 
        
        Consider the surjective homomorphism $\nu : \units[n] \to \units[m]$ defined in the lemma above. Since $(a, m) = 1$, we have $a \in \units[m]$, so there exists some $b \in \units[n]$ such that $\nu(b) = b \bmod m = a$. Then $\theta((h, k)) = (\sigma\inv(h), k^b)$ is a 2-sided inverse of $\psi$, so $\psi$ is bijective. \qh
    \end{enumerate}
\end{sol}

\begin{spex}[5.5.7]
    This exercise describes thirteen isomorphism types of groups of order 56. (It is not too difficult to show that every group of order 56 is isomorphic to one of these.)
    \begin{subexercises}
        \item Prove that there are three Abelian groups of order 56.
        \item Prove that every group of order 56 has either a normal Sylow 2-subgroup or a normal Sylow 7-subgroup.
        \item Construct the following non-Abelian groups of order 56 which have a normal Sylow 7-subgroup and whose Sylow 2-subgroup $S$ is as specified:
        \begin{itemize}
            \item one group when $S \cong Z_2 \times Z_2 \times Z_2$
            \item two non-isomorphic groups when $S \cong Z_4 \times Z_2$
            \item one group when $S \cong Z_8$
            \item two non-isomorphic groups when $S \cong Q_8$
            \item three non-isomorphic groups when $S \cong D_8$
        \end{itemize}
        [For a particular $S$, two groups are not isomorphic if the kernels of the maps from $S$ into $\aut(Z_7)$ are not isomorphic.]
        \item Let $G$ be a group of order 56 with a non-normal Sylow 7-subgroup. Prove that if $S$ is the Sylow 2-subgroup of $G$ then $S \cong Z_2 \times Z_2 \times Z_2$. [Let an element of order 7 act by conjugation on the seven non-identity elements of $S$ and deduce that they all have the same order.]
        \item Prove that there is a unique group of order 56 with a non-normal Sylow 7-subgroup. [For existence use the fact that $|GL_3(\f_2)| = 168$; for uniqueness use \exref{5.5.6}.]
    \end{subexercises}
\end{spex}

\begin{solenum}
    \item Note that $56 = 2^3 \cdot 7$. By the Fundamental Theorem of Finite Abelian Groups, the partitions of 3 are $3, 2 + 1, 1 + 1 + 1$, so the Abelian groups of order 56 are of the form $Z_8 \times Z_7$, $Z_4 \times Z_2 \times Z_7$, and $Z_2 \times Z_2 \times Z_2 \times Z_7$.
    \item Let $n_7$ be the number of Sylow 7-subgroups of $G$. Then $n_7 \equiv 1 \bmod 7$ and $n_7 \mid 2^3$, so $n_7 = 1$ or $8$. If $n_7 = 1$, then the unique Sylow 7-subgroup is normal. Suppose $n_7 = 8$. Then there are $8 \cdot 6 = 48$ elements of order 7 in $G$, so the remaining $56 - 48 = 8$ elements of $G$ must be the identity and the elements of a Sylow 2-subgroup. Hence, the Sylow 2-subgroup is unique, so it is normal.
    \item Let $P \in \syl_7(G)$ such that $P \nsub G$. Then $P = \gen p \cong Z_7$, and $\aut(P) \cong Z_6$.
    \begin{itemize}
        \item If $S \cong Z_2 \times Z_2 \times Z_2$, put $S = \gen a \times \gen b \times \gen c$. Let $\phi : S \to \aut(P)$ be a homomorphism. We disregard $a, b, c$ being sent to $\id$ as that is the trivial homomorphism. Since $a, b, c$ have order 2, they can only be sent to elements of order 1 or 2 in $\aut(P)$. Put $\alpha \in \aut(P)$ such that $|\alpha| = 2$, where $\alpha$ is unique since $\aut(P)$ is cyclic. Then we have the following non-trivial homomorphisms:
        \[\phi_1 = 
        \begin{cases}
            a \mapsto \alpha, \\
            b \mapsto \id, \\
            c \mapsto \id,
        \end{cases} \quad
        \phi_2 =
        \begin{cases}
            a \mapsto \alpha, \\
            b \mapsto \alpha, \\
            c \mapsto \id,
        \end{cases} \quad
        \phi_3 =
        \begin{cases}
            a \mapsto \alpha, \\
            b \mapsto \alpha, \\
            c \mapsto \alpha,
        \end{cases}\]
        Note that $\ker\phi_1 = \gen b \times \gen c$, $\ker\phi_2 = \gen{ab, c}$, and $\ker\phi_3 = \gen{ab, ac}$. Note that these are all isomorphic to $V_4$, hence the resulting semidirect products are isomorphic by \exref{5.5.6}. Without loss of generality, consider $\phi_1$. Then $a$ acts on $P$ by sending $p$ to the inversion automorphism of $P$, or $apa\inv = p\inv$. It is clear that $b$ and $c$ centralize $P$ and commute with $a$. Observe that $\gen{p, a \mid apa\inv = p\inv, a^2 = p^7 = 1} \cong D_{14}$ and has trivial intersection with $\gen b \times \gen c$, so
        \[G \cong P \semidir_{\phi_1} S \cong D_{14} \times Z_2 \times Z_2.\]
        \item Put $S = Z_4 \times Z_2 \cong \gen a \times \gen b$. Let $\alpha \in \aut(P)$ be the unique element of order 2 again. Then we have the following non-trivial homomorphisms:
        \[\phi_1 =
        \begin{cases}
            a \mapsto \alpha, \\
            b \mapsto \id,
        \end{cases} \quad
        \phi_2 = 
        \begin{cases}
            a \mapsto \id, \\
            b \mapsto \alpha,
        \end{cases} \quad
        \phi_3 =
        \begin{cases}
            a \mapsto \alpha, \\
            b \mapsto \alpha,
        \end{cases}\]
        We have $\ker\phi_1 = \gen{a^2} \times \gen b, \ker\phi_2 = \gen a$, and $\ker\phi_3 = \gen{ab}$, where $|ab| = 4$. Note that $\ker\phi_1 \cong V_4$, while $\ker\phi_2 \cong \ker\phi_3 \cong Z_4$. We may consider $\phi_2$ first. Then $b$ acts on $P$ by sending $p$ to the inversion automorphism of $P$, or $bpb\inv = p\inv$. It is clear that $a$ centralizes $P$ and commutes with $b$. Observe that $\gen{p, b \mid bpb\inv = p\inv, b^2 = p^7 = 1} \cong D_{14}$ and has trivial intersection with $\gen a$, so
        \[G \cong P \semidir_{\phi_2} S \cong D_{14} \times Z_4.\]
        For $\phi_1$, we see that $\gen b$ centralizes $P$, and $a$ acts on $P$ by sending $p$ to the inversion automorphism of $P$. Moreover, consider the subgroup $P_0 = \gen{p, a \mid apa\inv = p\inv, a^4 = p^7 = 1}$. It is easy to see that $P \cap \gen a = 1$, so $|P_0| = |P||\gen a| = 28$. Moreover, $pPp\inv = P$, and $aPa\inv = P$ since $apa\inv = p\inv$ under the action of $\phi_1$ restricted to $\gen a$. Hence, $P_0 \cong Z_7 \semidir Z_4$. It is clear that $P_0$ is normal in $G$, and $\gen b$ centralizes $P_0$, so
        \[G \cong P \semidir_{\phi_1} S \cong (Z_7 \semidir Z_4) \times Z_2.\]
        \item For $S \cong Z_8$, let $S = \gen s$. Since $\aut(P) \cong Z_6$, and $(8, 6) = 2$, then the only non-trivial homomorphism $\phi : S \to \aut(P)$ is the one that sends $s$ to the inversion automorphism $\alpha$ of $P$. Moreover, cyclicity of $S$ guarantees uniqueness of the homomorphism. Then $sps\inv = p\inv$, and
        \[G \cong P \semidir_\phi S \cong Z_7 \semidir Z_8.\]
        \item Let $S \cong Q_8 = \gen{i, j} = \set{1, i, i^2, i^3, j, ij, i^2j, i^3j}$. If we consider the trivial homomorphism, then $G \cong P \times S \cong Z_7 \times Q_8$. Now consider the non-trivial homomorphism that sends some generator of $S$ to $\alpha$. We have three cases:
        \[\phi_1 = 
        \begin{cases}
            i \mapsto \alpha, \\
            j \mapsto \id,
        \end{cases} \quad
        \phi_2 = 
        \begin{cases}
            i \mapsto \id, \\
            j \mapsto \alpha,
        \end{cases} \quad
        \phi_3 =
        \begin{cases}
            i \mapsto \alpha, \\
            j \mapsto \alpha,
        \end{cases}\]
        We have $\ker\phi_1 = \gen j, \ker\phi_2 = \gen i$, and $\ker\phi_3 = \gen{ij}$. These are all isomorphic to $Z_4$, so the semidirect products will be isomorphic to each other. Pick $\phi_1$ without loss of generality. Then $i$ acts on $P$, and $ipi\inv = p\inv$. Unlike the other cases, $j$ does not centralize $P$, since $jij\inv = i\inv$, so $jpj\inv = p\inv$. Hence, we have
        \[G \cong P \semidir_{\phi_1} S \cong Z_7 \semidir Q_8 = \set{p, i, j \mid i^4 = j^4 = 1, i^2 = j^2, iji\inv = j\inv, iji\inv = jpj\inv = p\inv}\]
        Note that non-triviality of the homomorphism implies that some $s \in S$ does not centralize $P$. If $G$ could be written as a direct product, then $S$ would centralize $P$, which is a contradiction. Hence, $G$ is not a direct product.
        \item Let $S \cong D_8 = \gen{r, s} = \set{1, r, r^2, r^3, s, rs, r^2s, r^3s}$. If we consider the trivial homomorphism, then $G \cong P \times S \cong Z_7 \times D_8$. Now consider the non-trivial homomorphism that sends some generator of $S$ to $\alpha$. We have three cases:
        \[\phi_1 = 
        \begin{cases}
            r \mapsto \alpha, \\
            s \mapsto \id,
        \end{cases} \quad
        \phi_2 =
        \begin{cases}
            r \mapsto \id, \\
            s \mapsto \alpha,
        \end{cases} \quad
        \phi_3 =
        \begin{cases}
            r \mapsto \alpha, \\
            s \mapsto \alpha,
        \end{cases}\]
        We have $\ker\phi_1 = \gen{r^2, s}, \ker\phi_2 = \gen r$, and $\ker\phi_3 = \gen{sr, r^2}$. Then $\ker\phi_1 \cong \ker\phi_3 \cong V_4$, while $\ker\phi_2 \cong Z_4$. Hence, the semidirect products corresponding to $\phi_1$ and $\phi_3$ are isomorphic to each other, but not to the semidirect product corresponding to $\phi_2$. Pick $\phi_1$ without loss of generality. Then $r$ acts on $P$, and $rpr\inv = p\inv$. Moreover, $s$ centralizes $P$, hence
        \[G \cong P \semidir_{\phi_1} S \cong Z_7 \semidir D_8 = \set{p, r, s \mid p^7 = r^4 = s^2 = 1, srs\inv = r\inv, rpr\inv = p\inv}.\]
        If we pick $\phi_2$, then $s$ acts on $P$, and $sps\inv = p\inv$. Moreover, note that $r$ centralizes $P$, hence $|rp| = 28$. Since $s(rp)s = srps\inv = (srs\inv)(sps\inv) = r\inv p\inv = (rp)\inv$, then
        \[G \cong P \semidir_{\phi_2} S \cong D_{56}.\]
    \end{itemize}
    \item Suppose $G$ has a non-normal Sylow 7-subgroup. We know that $S \in \syl_2(G)$ must be normal by part (b). Let $P \in \syl_7(G)$, and put $P = \gen p$. Let $p$ act by conjugation on the seven non-identity elements of $S$. Note that the stabilizer of any non-identity element $s \in S$ is a subgroup of $P$, hence $|G_s| = 1$ or 7. If $|G_s| = 7$, then the orbit is trivial, hence $psp\inv = s$, or $p = sps\inv \in P$, which contradicts that $|s|$ is a power of 2. Then $|G_s| = 1$, so the orbit $\oo_s$ of $s$ has size 7. Since $|S - 1| = 7$, then $\oo_s$ contains all non-identity elements of $S$. Hence, all non-identity elements of $S$ have the same order.
    
    If they all had order 8, then $s^2$ would have order 4, but $(s^2)^4 = s^8 = 1$. A similar reasoning eliminates the possibility of order 4. Hence, all non-identity elements of $S$ have order 2, so $S \cong Z_2 \times Z_2 \times Z_2$.
    \item Let $H = Z_2 \times Z_2 \times Z_2$ and $K = \gen k \cong Z_7$. Then $\aut(H) \cong \gl_3(\f_2)$, which has order 168. Let $\phi : K \to \aut(H)$ be a homomorphism. Since $|K| = 7$, then $|\phi(K)| = 1$ or 7. If $|\phi(K)| = 1$, then the resulting semidirect product is Abelian, which is a contradiction. Hence, $|\phi(K)| = 7$. Since $\aut(H)$ has order 168, then there exists a unique subgroup of order 7 in $\aut(H)$, call it $\psi$. Then $\phi : K \to \aut(H)$ given by $\phi(k) = \psi$ is a homomorphism, and the resulting semidirect product $H \semidir_\phi K$ has order $|H||K| = 8 \cdot 7 = 56$. We now claim that $K$ is not normal in $H \semidir_\phi K$.
    
    If $K \nsub G$, then $[s, k] \in K$ for all $s \in S$. Moreover, $[s, k] \in S$ since $S \nsub G$. Hence, $[s, k] \in S \cap K = 1$, so $[s, k] = 1$, or $sk = ks$ for all $s \in S$. Then $\phi(k) = \id$, which is a contradiction. Hence, $K$ is not normal in $H \semidir_\phi K$.

    To see that $H \semidir_\phi K$ is unique, let $\phi_1, \phi_2 : K \to \aut(H)$ be two homomorphisms such that $|\phi_1(K)| = |\phi_2(K)| = 7$. Then $\phi_1(K)$ and $\phi_2(K)$ are conjugate subgroups of $\aut(H)$. Since $|\aut(H) = 2^3 \cdot 3 \cdot 7$, then $\phi_1(K), \phi_2(K) \in \syl_7(\aut(H))$, so they are conjugate by the Sylow theorems. Hence, by \exref{5.5.6}, we have $H \semidir_{\phi_1} K \cong H \semidir_{\phi_2} K$. Therefore, $H \semidir_\phi K$ is unique.
\end{solenum}

\begin{exercise}[5.5.8]
    Construct a non-Abelian group of order 75. Classify all groups of order 75 (there are three of them). [Use \exref{5.5.6} to show that the non-Abelian group is unique.] (The classification of groups of order $pq^2$, where $p$ and $q$ are primes with $p < q$ and $p$ not dividing $q - 1$, is quite similar.)
\end{exercise}

\begin{sol}
    Let $P \in \syl_5(G)$ and $Q \in \syl_3(G)$. Then $|P| = 25$ and $|Q| = 3$. By the Sylow theorems, we have $n_5 \equiv 1 \bmod 5$ and $n_5 \mid 3$, so $n_5 = 1$ or 3. If $n_5 = 1$, then the unique Sylow 5-subgroup is normal. If $n_5 = 3$, then there are $3 \cdot 24 = 72$ elements of order 5 in $G$, so the remaining $75 - 72 = 3$ elements of $G$ must be the identity and the elements of a Sylow 3-subgroup. Hence, the Sylow 3-subgroup is unique, so it is normal. In either case, we have a normal Sylow subgroup, so $G \cong P \semidir Q$ or $G \cong Q \semidir P$. Moreover, $P \cong Z_{25}$ or $Z_5 \times Z_5$, and $Q \cong Z_3$. We then have two cases to consider:
    \begin{itemize}
        \item Suppose $P \cong Z_{25}$. Then $\aut(P) \cong \units[25]$, which has order 20. Since $|Q| = 3$, then the only homomorphism $\phi : Q \to \aut(P)$ is the trivial homomorphism, so $G \cong P \semidir Q \cong P \times Q \cong Z_{25} \times Z_3$.
        
        If we consider $G \cong Q \semidir P$, then we have a homomorphism $\phi : P \to \aut(Q)$. Since $|\aut(Q)| = 2$, then the only homomorphism $\phi : P \to \aut(Q)$ is the trivial homomorphism, so $G \cong Q \semidir P \cong Q \times P \cong Z_3 \times Z_{25}$.
        \item Suppose $P \cong Z_5 \times Z_5$ with $\aut(P) \cong \gl_2(\f_5)$, which has order 480. Since $|Q| = 3$, then there exists a non-trivial homomorphism $\phi : Q \to \aut(P)$, which is unique since $|\aut(P)| = 480$ and $(480, 3) = 3$. We may use similar reasoning as the previous exercise in that $\phi_1(Q)$ and $\phi_2(Q)$ are conjugate subgroups of $\aut(P)$, and $\phi_1(Q), \phi_2(Q) \in \syl_3(\aut(P))$, so they are conjugate by the Sylow theorems. Hence, by \exref{5.5.6}, we have $P \semidir_{\phi_1} Q \cong P \semidir_{\phi_2} Q$. Therefore, $P \semidir_\phi Q$ is unique, and $G \cong P \semidir_\phi Q$ is non-Abelian.
        
        If we consider $G \cong Q \semidir P$, then we have a homomorphism $\phi : P \to \aut(Q)$. Since $|\aut(Q)| = 2$, then the only homomorphism $\phi : P \to \aut(Q)$ is the trivial homomorphism, so $G \cong Q \semidir P \cong Q \times P \cong Z_3 \times (Z_5 \times Z_5)$.
    \end{itemize}
\end{sol}

\begin{exercise}[5.5.9]
    Show that the matrix
    \[
    \begin{bmatrix}
        0 & -1 \\
        1 & 4
    \end{bmatrix}\]
    is an element of order 5 in $\gl_2(\f_{19})$. Use this matrix to construct a non-Abelian group of order 1805, and give a presentation of this group. Classify groups of order 1805 (there are three isomorphism types). [Use \exref{5.5.6} to prove uniqueness of the non-Abelian group.] (A general method for finding elements of prime order in $\gl_n(\f_p)$ is described in the exercises in Section 12.2; this particular matrix of order 5 in $\gl_2(\f_{19})$ appears in \exref{12.2.16} of that section as an illustration of the method.)
\end{exercise}

\begin{sol}
    Let $M$ be the given matrix. It is easy to calculate that
    \[M^5 = 
    \begin{bmatrix}
       -56 & -209 \\
       209 & 780
    \end{bmatrix},\]
    where, once entries are taken mod 19, results in $I_2$, hence $|M| = 5$. Moreover, we know that $1805 = 19^2 \cdot 5$. Since $n_{19} \equiv 1 \bmod 19$ and $n_{19} \mid 5$, then $n_{19} = 1$, so $P \in \syl_{19}(G)$ is normal. Moreover, we have $P \cong Z_{361}$ or $P \cong Z_{19} \times Z_{19}$. If $P \cong Z_{361}$, then $\aut(P) \cong \units[361]$, which has order 19 * 18 = 342. Since $|Q| = 5$, then the only homomorphism $\phi : Q \to \aut(P)$ is the trivial homomorphism, so $G \cong P \semidir Q \cong P \times Q \cong Z_{361} \times Z_5$.

    If $P \cong Z_{19} \times Z_{19}$, then $\aut(P) \cong \gl_2(\f_{19})$, which has order 123120. If we take $\phi$ to be the trivial homomorphism, then $G \cong Z_{19} \times Z_{19} \times Z_5$. Now, since $|Q| = 5$, then there exists a non-trivial homomorphism $\phi : Q \to \aut(P)$, which is unique since $|\aut(P)| = 123120$ and $(123120, 5) = 5$. We may use similar reasoning as the previous exercise in that $\phi_1(Q)$ and $\phi_2(Q)$ are conjugate subgroups of $\aut(P)$, and $\phi_1(Q), \phi_2(Q) \in \syl_5(\aut(P))$, so they are conjugate by the Sylow theorems. Hence, by \exref{5.5.6}, we have $P \semidir_{\phi_1} Q \cong P \semidir_{\phi_2} Q$. Therefore, $P \semidir_\phi Q$ is unique, and $G \cong P \semidir_\phi Q$ is non-Abelian.
\end{sol}

\begin{spex}[5.5.10]
    This exercise classifies the groups of order 147 (there are six isomorphism types).
    \begin{subexercises}
        \item Prove that there are two Abelian groups of order 147.
        \item Prove that every group of order 147 has a normal Sylow 7-subgroup.
        \item Prove that there is a unique non-Abelian group of order 147 whose Sylow 7-subgroup is cyclic.
        \item Let
        \[t_1 = 
        \begin{bmatrix}
            2 & 0 \\
            0 & 1
        \end{bmatrix}, \quad
        t_2 = 
        \begin{bmatrix}
            1 & 0 \\
            0 & 2
        \end{bmatrix}\]
        be elements of $\gl_2(\f_7)$. Prove $P = \gen{t_1, t_2}$ is a Sylow 3-subgroup of $\gl_2(\f_7)$, and that $P \cong Z_3 \times Z_3$. Deduce that every subgroup of order 3 is conjugate in $\gl_2(\f_7)$ to a subgroup of $P$.
        \item By Example 3 in Section 1, the group $P$ has four subgroups of order 3, and these are: $P_1 = \gen{t_1}, P_2 = \gen{t_2}, P_3 = \gen{t_1t_2}$, and $P_4 = \gen{t_1t_2^2}$. For $i = 1, 2, 3, 4$, let $G_i = (Z_7 \times Z_7) \semidir_{\phi_i} P_i$, where $\phi_i$ is an isomorphism of $Z_3$ with the subgroup $P_i$ of $\aut(Z_7 \times Z_7)$. For each $i$, describe $G_i$ in terms of generators and relations. Deduce that $G_1 \cong G_2$.
        \item Prove that $G_1$ is not isomorphic to either $G_3$ or $G_4$. [Show that the center of $G_1$ has order 7, whereas the centers of $G_3$ and $G_4$ are trivial.]
        \item Prove that $G_3$ is not isomorphic to $G_4$.
        \item Classify the groups of order 147 by showing that the six non-isomorphic groups described above (two from part (a), one from part (c), and $G_1, G_3$, and $G_4$) are all the groups of order 147. [Use \exref{5.5.6} and part (d).] (The classification of groups of order $pq^2$, where $p$ and $q$ are primes with $p < q$ and $p \mid q - 1$, is quite similar.)
    \end{subexercises}
\end{spex}

\begin{solenum}
    \item Note that $147 = 3 \cdot 7^2$. By the Fundamental Theorem of Finite Abelian Groups, the partitions of 2 are $2, 1 + 1$, so the Abelian groups of order 147 are of the form $Z_{49} \times Z_3$ and $Z_7 \times Z_7 \times Z_3$.
    \item Let $n_7$ be the number of Sylow 7-subgroups of $G$. Then $n_7 \equiv 1 \bmod 7$ and $n_7 \mid 3$, so $n_7 = 1$. Hence, the unique Sylow 7-subgroup is normal.
    \item Let $P \cong Z_{49}$ and $Q \cong Z_3$. Then $\aut(P) \cong \units[49]$, which has order $7 \cdot 6 = 42$. Since $|Q| = 3$, then there exists non-trivial homomorphisms $\phi : Q \to \aut(P)$, where $q \in Q = \gen q$ is mapped to a generator of the unique subgroup of order 3 in $\aut(P)$, say $\alpha$. Then either $q \mapsto \alpha$ or $q \mapsto \alpha^2$. Let $\phi_1, \phi_2 : Q \to \aut(P)$ be the two homomorphisms defined by $\phi_1(q) = \alpha$ and $\phi_2(q) = \alpha^2$. Since $\phi_1(Q) = \phi_2(Q) = \gen\alpha$, then by \exref{5.5.6}, we have $P \semidir_{\phi_1} Q \cong P \semidir_{\phi_2} Q$. Hence, there is a unique non-Abelian group of order 147 whose Sylow 7-subgroup is cyclic.
    \item Recall that $|\gl_2(\f_7)| = (7^2 - 1)(7^2 - 7) = 48 \cdot 42 = 2016 = 2^5 \cdot 3^2 \cdot 7$. Observe that
    \[t_1^3 = 
    \begin{bmatrix}
        8 & 0 \\
        0 & 1
    \end{bmatrix}, \quad 
    t_2^3 = 
    \begin{bmatrix}
        1 & 0 \\
        0 & 8
    \end{bmatrix}, \quad
    t_1t_2 = 
    \begin{bmatrix}
        2 & 0 \\
        0 & 2
    \end{bmatrix} = t_2t_1.\]
    Taking entries mod 7, then $t_1^3 = t_2^3 = I_2$, so $|t_1| = |t_2| = 3$. Moreover, $t_1t_2 = t_2t_1$, so $P = \gen{t_1}\gen{t_2}$. To see that this is isomorphic to $\gen{t_1} \times \gen{t_2}$, note that $\gen{t_1} \cap \gen{t_2} = 1$. Moreover, commutativity shows that ${t_1} \nsub P$ and ${t_2} \nsub P$, so $P \cong \gen{t_1} \times \gen{t_2} \cong Z_3 \times Z_3$. Since $|P| = 9$, then $P$ is a Sylow 3-subgroup of $\gl_2(\f_7)$.

    Let $R$ be a subgroup of order 3 in $\gl_2(\f_7)$. Then $R \leq gPg\inv$ for some $g \in \gl_2(\f_7)$ by the Sylow theorems. Then $gRg\inv \leq P$, so $gRg\inv$ is a subgroup of order 3 in $P$. Hence, every subgroup of order 3 is conjugate in $\gl_2(\f_7)$ to a subgroup of $P$.
    \item Let $Z_7 \times Z_7 = \gen a \times \gen b$. Then $\aut(Z_7 \times Z_7) \cong \gl_2(\f_7)$, and $P_i$ is a subgroup of $\aut(Z_7 \times Z_7)$ for $i = 1, 2, 3, 4$. Let $\phi_i : Z_3 \to P_i$ be an isomorphism. We check each $i$:
    \begin{itemize}
        \item For $i = 1$, we have $\phi_1 : Z_3 \to P_1$ given by $\phi_1(q) = t_1$. Then $qaq\inv = \phi_1(q)(a) = t_1(a) = a^2$, and $qbq\inv = \phi_1(q)(b) = t_1(b) = b$. Hence, we have for a na\"ive presentation of $G_1$:
        \[G_1 = \set{a, b, q \mid a^7 = b^7 = q^3 = 1, ab = ba, qaq\inv = a^2, qbq\inv = b}.\]
        However, note that $b$ is central in $G_1$ as it commutes with both $a$ and $q$. It follows that $\gen b \nsub G_1$ with $|\gen b| = 7$. Observe that $\gen b \cap \gen{a, q} = 1$ by construction. To see that $\gen{a, q} \nsub G_1$, observe that $b\gen{a, q}b\inv = \gen{a, q}$ because $b \in Z(G_1)$. Then $G_1 \cong Z_7 \times \gen{a, q}$.

        If we consider a homomorphism $\psi : Z_3 \to \aut(Z_7)$, then $q \mapsto \alpha$ or $q \mapsto \alpha^2$ such that $\alpha$ generates the unique subgroup of order 3 in $\aut(Z_7)$. In either case, we have $\gen{a, q} \cong Z_7 \semidir_\psi Z_3$ as it satisfies the relations $qaq\inv = a^2$ in the na\"ive presentation. Hence, we have
        \[G_1 \cong Z_7 \times (Z_7 \semidir_\psi Z_3).\]
        \item For $i = 2$, simply note that $t_2$ takes the place of $t_1$ in the previous case, $aq = qa$, and $qbq\inv = b^2$. Hence, we have
        \[G_2 = \set{a, b, q \mid a^7 = b^7 = q^3 = 1, ab = ba, qaq\inv = a, qbq\inv = b^2}.\]
        By similar reasoning as the previous case, we have $G_2 \cong Z_7 \times (Z_7 \semidir_\psi Z_3)$, so $G_1 \cong G_2$.
        \item For $i = 3$, we have $\phi_3 : Z_3 \to P_3$ given by $\phi_3(q) = t_1t_2$. Then $qaq\inv = \phi_3(q)(a) = t_1t_2(a) = t_1(a) = a^2$, and $qbq\inv = \phi_3(q)(b) = t_1t_2(b) = t_1(b^2) = b^2$. A presentation of $G_3$ is
        \[G_3 = \set{a, b, q \mid a^7 = b^7 = q^3 = 1, ab = ba, qaq\inv = a^2, qbq\inv = b^2}.\]
        \item For $i = 4$, we have $\phi_4 : Z_3 \to P_4$ given by $\phi_4(q) = t_1t_2^2$. Then $qaq\inv = \phi_4(q)(a) = t_1t_2^2(a) = t_1(a) = a^2$, and $qbq\inv = \phi_4(q)(b) = t_1t_2^2(b) = t_1(b^4) = b^4$. Then a presentation of $G_4$ is
        \[G_4 = \set{a, b, q \mid a^7 = b^7 = q^3 = 1, ab = ba, qaq\inv = a^2, qbq\inv = b^4}.\]
    \end{itemize}
    \item The prior discussion showed that $Z(G_1) = \gen b$, so $|Z(G_1)| = 7$. Now, consider $G_3$, and let $z \in Z(G_3)$. We first consider when $z \in Z(G_3) \cap (\gen a \times \gen b)$. Then $z = a^ib^j$ for some $0 \leq i, j \leq 6$. Since $z$ commutes with $q$, then 
    \[a^ib^j = z = qzq\inv = qa^ib^jq\inv = qaq\inv qbq\inv = a^{2i}b^{2j}.\]
    Then this implies that $2i \equiv i \bmod 7$ and $2j \equiv j \bmod 7$, so $i = j =0$. Hence, $Z(G_3) \cap (\gen a \times \gen b) = 1$. Now, consider when $z \in Z(G_3) - (\gen a \times \gen b)$. Then $z = a^ib^jq^k$ for some $0 \leq i, j \leq 6$ and $1 \leq k \leq 2$. Since $z$ commutes with $a$, then 
    \[a^ib^jq^k = z = aza\inv = aa^ib^jq^ka\inv = a^{i+1}b^jq^ka\inv = a^{i+1}b^jq^ka^{-1}.\]
    However, since $q$ does not centralize $\gen a$, then this is a contradiction. Hence, $Z(G_3) - (\gen a \times \gen b) = 1$, so $Z(G_3) = 1$. A similar argument shows that $Z(G_4) = 1$.
    \item To show that $G_3 \not\cong G_4$, we first observe that $G_3$ and $G_4$ have the same subgroups of order 7, since $Z_7 \times Z_7$ is the unique Sylow 7-subgroup of both groups. Hence, $G_3$ and $G_4$ have the list of subgroups
    \[\ss = \set{\gen a, \gen b, \gen{ab}, \gen{ab^2}, \gen{ab^3}, \gen{ab^4}, \gen{ab^5}, \gen{ab^6}}.\]
    Let $q$ act on $\ss$ by conjugation. In $G_3$, for any $\gen{ab^i} \in \ss$, then
    \[q(ab^i)q\inv = (qaq\inv)(qb^iq\inv) = a^2b^{2i}.\]
    It is clear that $a^2b^{2i} \in \gen{ab^i}$, so $q$ stabilizes $\gen{ab^i}$. Hence, $q$ stabilizes all subgroups of order 7 in $G_3$, and each subgroup of order 7 is normal in $G_3$. In $G_4$, consider $\gen{ab} \in \ss$. Then
    \[q(ab)q\inv = (qaq\inv)(qbq\inv) = a^2b^4 \notin \gen{ab},\]
    so $q$ does not stabilize $\gen{ab}$. Hence, $\gen{ab}$ is not normal in $G_4$. Since $G_3$ and $G_4$ have different numbers of normal subgroups of order 7, then $G_3 \not\cong G_4$.
    \item By part (b), every group of order 147 has a normal Sylow 7-subgroup. Let $P \in \syl_7(G)$ and $Q \in \syl_3(G)$. Since $P \cap Q = 1$ because (49, 3) = 1, and $|G| = |P||Q|$, then $G = PQ$, and $G \cong P \semidir Q$. If the homomorphism is trivial, then we have the two Abelian groups from part (a). If the homomorphism is non-trivial and $P$ is cyclic, then we have the unique non-Abelian group from part (c). If the homomorphism is non-trivial and $P$ is not cyclic, then $\phi(Q)$ is a subgroup of order 3 in $\aut(P) \cong \gl_2(\f_7)$. By part (d), every subgroup of order 3 is conjugate in $\gl_2(\f_7)$ to a subgroup of $P$, so by \exref{5.5.6} we have the three non-Abelian groups $G_1, G_3$, and $G_4$ from part (e). Hence, the six non-isomorphic groups described above are all the groups of order 147.
\end{solenum}

\begin{exercise}[5.5.11]
    Classify groups of order 28 (there are four isomorphism types).
\end{exercise}

\begin{sol}
    Note $28 = 2^2 \cdot 7$. We have the Abelian groups $Z_7 \times Z_4$ and $Z_7 \times Z_2 \times Z_2$. Let $P \in \syl_7(G)$ and $S \in \syl_2(G)$. Then $|P| = 7$ and $|S| = 4$. By the Sylow theorems, we have $n_7 \equiv 1 \bmod 7$ and $n_7 \mid 4$, so $n_7 = 1$. Hence, the unique Sylow 7-subgroup is normal. Moreover, we have $P \cong Z_7$ and $S \cong Z_4$ or $Z_2 \times Z_2$. We then have two cases to consider:
    \begin{itemize}
        \item Suppose $S \cong Z_4 = \gen s$. If we consider the trivial homomorphism, then $G \cong P \times S \cong Z_7 \times Z_4$. Now consider the non-trivial homomorphism that sends $s$ to $\alpha$, a generator of the unique subgroup of order 2 in $\aut(P)$. Then $s$ acts on $P$, and $sps\inv = p\inv$. Hence, we have
        \[G \cong P \semidir S \cong Z_7 \semidir Z_4 = \set{p, s \mid p^7 = s^4 = 1, sps\inv = p\inv}.\]
        \item Suppose $S \cong Z_2 \times Z_2 = \gen s \times \gen t$. If we consider the trivial homomorphism, then $G \cong P \times S \cong Z_7 \times (Z_2 \times Z_2)$. Now consider a non-trivial homomorphism that sends some generator of $S$ to $\alpha$. We have the three homomorphisms
        \[\theta_1 = 
        \begin{cases}
            s \mapsto \alpha, \\
            t \mapsto \id,
        \end{cases} \quad
        \theta_2 = 
        \begin{cases}
            s \mapsto \id, \\
            t \mapsto \alpha,
        \end{cases} \quad
        \theta_3 = 
        \begin{cases}
            s \mapsto \alpha, \\
            t \mapsto \alpha,
        \end{cases}\]
        Observe that $\ker\theta_1 = \gen t, \ker\theta_2 = \gen s$, and $\ker\theta_3 = \gen{st}$. Hence, the semidirect products corresponding to $\theta_1, \theta_2$, and $\theta_3$ are all isomorphic to each other. Pick $\theta_1$ without loss of generality. Then $s$ acts on $P$, and $sps\inv = p\inv$. Moreover, $t$ centralizes $P$, hence $|tp| = 14$. Since
        \[s(tp)s\inv = stps\inv = (sts\inv)(sps\inv) = tp\inv = (tp)\inv,\]
        then
        \[G \cong P \semidir_{\theta_1} S \cong \set{p, s, t \mid p^7 = s^2 = t^2 = 1, st = ts, sps\inv = p\inv, tpt\inv = p} \cong D_{28}. \qh\] 
    \end{itemize}
\end{sol}

\begin{exercise}[5.5.12]
    Classify the groups of order 20 (there are five isomorphism types).
\end{exercise}

\begin{sol}
    We have the Abelian groups $Z_5 \times Z_4$ and $Z_5 \times Z_2 \times Z_2$. Let $P \in \syl_5(G)$ and $S \in \syl_2(G)$. Then $|P| = 5$ and $|S| = 4$. By the Sylow theorems, we have $n_5 \equiv 1 \bmod 5$ and $n_5 \mid 4$, so $n_5 = 1$. Hence, the unique Sylow 5-subgroup is normal. Moreover, we have $P \cong Z_5$ and $S \cong Z_4$ or $Z_2 \times Z_2$. We then have two cases to consider:
    \begin{itemize}
        \item If $S \cong Z_4 = \gen s$, then $s$ has four candidates to be sent to: $\id, \alpha$, $\alpha^2$, or $\alpha^3$, where $\aut(Z_5) = \gen\alpha$. Define the homomorphisms
        \[\phi_1 : s \mapsto \id, \quad \phi_2 : s \mapsto \alpha, \quad \phi_3 : s \mapsto \alpha^2, \quad \phi_4 : s \mapsto \alpha^3\]
        Since $\phi_1$ is trivial, then $G \cong P \semidir_{\phi_1} S \cong P \times S \cong Z_5 \times Z_4$. Also observe that $\ker\phi_2 = \ker\phi_4 = 1$ while $\ker\phi_3 = \gen{s^2}$. Then the resulting semidirect products from $\phi_2$ and $\phi_4$ are isomorphic, since $\phi_2(S) = \phi_4(S) = \gen\alpha$, so by \exref{5.5.6}, we have $P \semidir_{\phi_2} S \cong P \semidir_{\phi_4} S$. Moreover, $\ker\phi_3 = \gen{s^2}$, which is the unique subgroup of order 2 in $\aut(P)$, so $P \semidir_{\phi_3} S$ is not isomorphic to $P \semidir_{\phi_2} S$. Hence, we have two non-Abelian groups of order 20, and we have
        \[G \cong P \semidir_{\phi_2} S \cong \set{p, s \mid p^5 = s^4 = 1, sps\inv = p^2} \cong F_{20}.\]
        In the case of $\phi_3$, we have $sps\inv = p\inv$. Then
        \[G \cong P \semidir_{\phi_3} S \cong \set{p, s \mid p^5 = s^4 = 1, sps\inv = p\inv} \cong Z_5 \semidir Z_4.\]
        \item If $S \cong Z_2 \times Z_2 = \gen s \times \gen t$, then we have the trivial homomorphism $\theta_1 : S \to \aut(P)$, which gives $G \cong P \semidir_{\theta_1} S \cong P \times S \cong Z_5 \times (Z_2 \times Z_2)$. If $\alpha^2$ is the generator of the unique subgroup of order 2 in $\aut(P)$, then we have the three non-trivial homomorphisms
        \[\theta_2 = 
        \begin{cases}
            s \mapsto \alpha^2, \\
            t \mapsto \id,
        \end{cases} \quad
        \theta_3 = 
        \begin{cases}
            s \mapsto \id, \\
            t \mapsto \alpha^2,
        \end{cases} \quad
        \theta_4 =
        \begin{cases}
            s \mapsto \alpha^2, \\
            t \mapsto \alpha^2,
        \end{cases}\]
        This is the same situation as in the previous exercise, since $\ker\theta_2 = \gen t, \ker\theta_3 = \gen s$, and $\ker\theta_4 = \gen{st}$. Hence, the semidirect products corresponding to $\theta_2, \theta_3$, and $\theta_4$ are all isomorphic to each other. Pick $\theta_2$ without loss of generality. Then $s$ acts on $P$, and $sps\inv = p\inv$. Moreover, $t$ centralizes $P$, hence $|tp| = 10$. Since
        \[s(tp)s\inv = stps\inv = (sts\inv)(sps\inv) = tp\inv = (tp)\inv,\]
        then
        \[G \cong P \semidir_{\theta_2} S \cong \set{p, s, t \mid p^5 = s^2 = t^2 = 1, st = ts, sps\inv = p\inv, tpt\inv = p} \cong D_{20}. \qh\]
    \end{itemize}
\end{sol}

\begin{exercise}[5.5.13]
    Classify groups of order $4p$, where $p$ is a prime greater than 3. [There are four isomorphism types when $p \equiv 3 \bmod 4$, and five isomorphism types when $p \equiv 1 \bmod 4$.]
\end{exercise}

\begin{sol}
    We have a framework for each case of $p$, namely \exref{5.5.11} when $p \equiv 3 \bmod 4$, and \exref{5.5.12} when $p \equiv 1 \bmod 4$. The only difference is that we have $P \in \syl_p(G)$ and $S \in \syl_2(G)$, where $|P| = p$ and $|S| = 4$. By the Sylow theorems, we have $n_p \equiv 1 \bmod p$ and $n_p \mid 4$, so $n_p = 1$. Hence, the unique Sylow $p$-subgroup is normal. Moreover, we have $P = \gen q \cong Z_p$ and $S \cong Z_4$ or $Z_2 \times Z_2$. We then have two cases to consider:
    \begin{itemize}
        \item If $S \cong Z_4 = \gen s$, then $s$ has different candidates to be sent to in $\aut(P) = \units[p]$, which has order $p - 1$. If $p \equiv 3 \bmod 4$, then $p - 1 \equiv 2 \bmod 4$, hence $\phi(s)$ can only be sent to $\id$ or the unique element of order 2 in $\aut(P)$. If we have the latter case, then $sqs\inv = q\inv$, and we have
        \[G \cong P \semidir S \cong Z_p \semidir Z_4 = \set{q, s \mid q^p = s^4 = 1, sqs\inv = q\inv}.\]
        
        If $p \equiv 1 \bmod 4$, then $p - 1 \equiv 0 \bmod 4$, hence $\phi(s)$ can be sent to $\id$ or one of the two elements of order 4 in $\aut(P)$ (which are isomorphic since their images generate the same subgroup), or the unique element of order 2 in $\aut(P)$. If $\aut(P) = \gen\alpha$, let $\rho = \alpha^{(p - 1)/4}$, which has order 4. Then we have the three non-trivial homomorphisms
        \[\phi_1 : s \mapsto \rho, \quad \phi_2 : s \mapsto \rho^3, \quad \phi_3 : s \mapsto \rho^2.\]
        We know that $\phi_1$ and $\phi_2$ give isomorphic semidirect products, since $\phi_1(S) = \phi_2(S) = \gen\rho$. Hence, we have two non-Abelian groups of order $4p$, and we have
        \[G \cong P \semidir_{\phi_1} S \cong \set{q, s \mid q^p = s^4 = 1, sqs\inv = q^k}.\]
        for some $k \in \units[p]$ such that $|k| = 4$. In the case of $\phi_3$, we have $sqs\inv = q\inv$. Then
        \[G \cong P \semidir_{\phi_3} S \cong \set{q, s \mid q^p = s^4 = 1, sqs\inv = q\inv}.\]
        \item If $S \cong Z_2 \times Z_2 = \gen s \times \gen t$, then we have the trivial homomorphism $\theta_1 : S \to \aut(P)$, which gives $G \cong P \semidir_{\theta_1} S \cong P \times S \cong Z_p \times (Z_2 \times Z_2)$. If $\rho^2$ is the generator of the unique subgroup of order 2 in $\aut(P)$, then we have the three non-trivial homomorphisms
        \[\theta_2 = 
        \begin{cases}
            s \mapsto \rho^2, \\
            t \mapsto \id,
        \end{cases} \quad
        \theta_3 =
        \begin{cases}
            s \mapsto \id, \\
            t \mapsto \rho^2,
        \end{cases} \quad
        \theta_4 =
        \begin{cases}
            s \mapsto \rho^2, \\
            t \mapsto \rho^2,
        \end{cases}\]
        Again, we know that $\ker\theta_2 = \gen t, \ker\theta_3 = \gen s$, and $\ker\theta_4 = \gen{st}$. Hence, the semidirect products corresponding to $\theta_2, \theta_3$, and $\theta_4$ are all isomorphic to each other. Pick $\theta_2$ without loss of generality. Then $s$ acts on $P$, and $sqs\inv = q\inv$. Moreover, $t$ centralizes $P$, hence $|tq| = 2p$. Since
        \[s(tq)s\inv = stqs\inv = (sts\inv)(sqs\inv) = tq\inv = (tq)\inv,\]
        then
        \[G \cong P \semidir_{\theta_2} S \cong \set{q, s, t \mid (tq)^2 = s^2 = 1, st = ts, sqs\inv = q\inv, tqt\inv = q} \cong D_{4p}. \qh\]
    \end{itemize}
\end{sol}

\begin{spex}[5.5.14]
    This exercise classifies the groups of order 60 (there are thirteen isomorphism types). Let $G$ be a group of order 60, let $P$ be a Sylow 5-subgroup of $G$, and let $Q$ be a Sylow 3-subgroup of $G$.
    \begin{subexercises}
        \item Prove that if $P$ is not normal in $G$, then $G \cong A_5$. [See Section 4.5.]
        \item Prove that if $P \nsub G$ but $Q$ is not normal in $G$, then $G \cong A_4 \times Z_5$. [Show in this case that $P \leq Z(G), G/P \cong A_4$, a Sylow 2-subgroup $T$ of $G$ is normal, and $TQ \cong A_4$.]
        \item Prove that if both $P$ and $Q$ are normal in $G$, then $G \cong Z_{15} \semidir T$, where $T \cong Z_4$ or $Z_2 \times Z_2$. Show in this case that there are six isomorphism types when $T$ is cyclic (one Abelian), and there are five isomorphism types when $T$ is the Klein 4-group (one Abelian). [Use the same ideas as in the classification of groups of order 30 and 20.]
    \end{subexercises}
\end{spex}

\begin{solenum}
    \item Clear, from Corollary 4.22.
    \item If $Q$ is not normal in $G$, then $n_3 \equiv 1 \bmod 3$ and $n_3 \mid 20$ implies that $n_3 = 4$ or 10. If $n_3 = 10$, then since $P \nsub G$, we know that $PQ_i$ is a subgroup of $G$ that is cyclic. But then $PQ_i \cap PQ_j = 5$ intersect at $P$, which means there would be $10 \cdot 10$ elements in $G$, which would be a contradiction. Hence, $n_3 = 4$. Again, each $PQ_i$ intersect non-trivially at $P$, and $P$ commutes with the non-identity elements of each $PQ_i$. Then there are $5 + 10 \cdot 4 = 45$ elements in $C_G(P)$. Since $C_G(P) \leq G$, then $|C_G(P)| = 60$ so that $C_G(P) = G$, showing that $P \leq Z(G)$.
    
    Consider $G/P$, where $|G/P| = 12$. If we consider $\pi$ to be the natural projection map from $G$ to $G/P$, we claim that the map $\psi : \syl_3(G) \to \syl_3(G/P)$ given by $\psi(Q_i) \to \pi(Q_i)$ is a bijection. Note that this is a well-defined map, since $\pi(Q_i) = Q_iP/P$ is a subgroup of $G/P$, and since $Q_iP \leq G$ with order 15, then $Q_iP/P$ has order 3, so $\pi(Q_i) \in \syl_3(G/P)$. To see that $\psi$ is injective, suppose that $\psi(Q_i) = \psi(Q_j)$, then $Q_iP/P = Q_jP/P$, which implies that $Q_iP = Q_jP$. Since $|Q_iP| = |Q_jP| = 15$, and $n_3(Q_iP) \equiv 1 \bmod 3$ and $n_3(Q_iP) \mid 5$, then $n_3(Q_iP) = 1$. Hence, $Q_i = Q_j$, so $\psi$ is injective. Since $n_3(G) = 4$, then $|\syl_3(G/P)| \geq 4$. By injectivity, it must be that $\psi$ is a bijection, hence $n_3(G/P) = 4$. We now consider the classification of groups of order 12, which are $Z_{12}, Z_2 \times Z_6, D_{12}, A_4$, and $Z_3 \semidir Z_4$.
    \begin{itemize}
        \item We know $Z_{12}$ and $Z_2 \times Z_6$ are Abelian, so any group of order 3 in these groups is normal, which is a contradiction.
        \item For $D_{12} \cong S_3 \times Z_2$, let $\alpha \in S_3$ be a 3-cycle. Then $\gen{(\alpha, 1)}$ is a subgroup of order 3. Recall that conjugating $\alpha$ by any cycle in $S_3$ gives $\alpha$ or $\alpha^2$, hence $\gen{(\alpha, 1)}$ is normal in $D_{12}$, which is a contradiction.
        \item For $Z_3 \semidir Z_4$, we know that $Z_3$ is normal in this group, which is a contradiction.
    \end{itemize}
    Then the only possibility is that $G/P \cong A_4$.

    Let $T_1, T_2 \in \syl_2(G)$ be distinct, and consider $T_1P$ and $T_2P$. Since $P \leq Z(G)$, then $T_iP \cong T_i \times P$ for $i = 1, 2$. If there was some $t \in T_1$ such that $t \in T_2P$, observe that $|t|$ is divisible by 4. Since all non-identity elements of $P$ have order 5, then $t \in T_2$. But $T_1$ and $T_2$ are distinct by assumption, hence there must be some $t^* \in T_1$ such that $t^* \notin T_2P$. Then $T_1P \neq T_2P$, and $T_1P/P, T_2P/P \in \syl_2(G/P)$. Since $n_2(G/P) = 1$, then $T_1P/P = T_2P/P$, which implies that $T_1P = T_2P$. But this is a contradiction, so there is a unique Sylow 2-subgroup of $G$, which we call $T$.

    Consider $TQ$. It is clear that $T \cap Q = 1$, and $|TQ| = 12$. If $Q \nsub TQ$, then since $P \leq Z(G)$ also normalizes $Q$, then $Q \nsub G$, which is a contradiction. Then $n_3(TQ) \neq 1$, and by a similar argument as above, we conclude that $TQ \cong A_4$. If we now consider $\phi : TQ \to \aut(P)$, note that $\phi$ is trivial since $TQ$ is generated by 3-cycles, while $\aut(P) \cong Z_4$ has no elements of order 3. Then $G \cong P \times TQ \cong Z_5 \times A_4$.
    \item If $P$ and $Q$ are both normal in $G$, then we have $PQ \nsub G$. Since $P \cap Q = 1$, then $PQ \cong P \times Q \cong Z_5 \times Z_3 \cong Z_{15}$. Put $P = \gen p$ and $Q = \gen q$. Let $T \in \syl_2(G)$, and consider a homomorphism from $T$ to $\aut(Z_{15})$, where $\aut(Z_{15}) \cong Z_4 \times Z_2$. From the text, we know that there are exactly 3 automorphisms of order 2:
    \[\alpha_1 = 
    \begin{cases}
        p \mapsto p\inv \\
        q \mapsto q
    \end{cases} \quad
    \alpha_2 = 
    \begin{cases}
        p \mapsto p \\
        q \mapsto q\inv
    \end{cases} \quad
    \alpha_3 = 
    \begin{cases}
        p \mapsto p\inv \\
        q \mapsto q\inv
    \end{cases}\]
    Moreover, there are 4 automorphisms of order 4:
    \[\beta_1 = 
    \begin{cases}
        p \mapsto p^2 \\
        q \mapsto q
    \end{cases} \quad
    \beta_2 =
    \begin{cases}
        p \mapsto p^2 \\
        q \mapsto q\inv
    \end{cases} \quad
    \beta_3 =
    \begin{cases}
        p \mapsto p^3 \\
        q \mapsto q
    \end{cases} \quad
    \beta_4 =
    \begin{cases}
        p \mapsto p^3 \\
        q \mapsto q\inv
    \end{cases}\]
    However, note that $\gen{\beta_1} = \gen{\beta_3}$ and $\gen{\beta_2} = \gen{\beta_4}$, so there are only two distinct subgroups of order 4 in $\aut(Z_{15})$. Without loss of generality, we may consider the automorphisms $\alpha_1, \alpha_2, \alpha_3, \beta_1$, and $\beta_2$. We then have the following cases to consider for the Sylow 2-subgroup $T$:
    \begin{itemize}
        \item Suppose $T = \gen t \cong Z_4$. Then we have the trivial homomorphism, which gives $G \cong Z_{15} \times Z_4$. Then we have the non-trivial homomorphisms
        \[\phi_1 : t \mapsto \alpha_1, \quad \phi_2 : t \mapsto \alpha_2, \quad \phi_3 : t \mapsto \alpha_3, \quad \phi_4 : t \mapsto \beta_1, \quad \phi_5 : t \mapsto \beta_2\]
        \begin{enumerate}
            \item For $G_1 = Z_{15} \semidir_{\phi_1} Z_4$, we have $tpt\inv = p\inv$ and $tqt\inv = q$. Since $q \in G_1$ is centralized, we have the factorization that $\gen{p, t} \cong Z_5 \semidir Z_4$. However, note that $\aut(Z_5) \cong Z_4$, so we have the presentation
            \[G_1 \cong \gen{p, t} \times \gen q \cong \hol(Z_5) \times Z_3.\]
            \item For $G_2 = Z_{15} \semidir_{\phi_2} Z_4$, we have $tpt\inv = p$ and $tqt\inv = q\inv$. Since $p \in G_2$ is centralized, we have the factorization
            \[G_2 \cong \gen{q, t} \times \gen p \cong (Z_3 \semidir Z_4) \times Z_5.\]
            \item For $G_3 = Z_{15} \semidir_{\phi_3} Z_4$, we have $tpt\inv = p\inv$ and $tqt\inv = q\inv$. Then
            \[G_3 \cong Z_{15} \semidir Z_4 = \set{p, q, t \mid p^5 = q^3 = t^4 = 1, pq = qp, tpt\inv = p\inv, tqt\inv = q\inv}.\]
            \item For $G_4 = Z_{15} \semidir_{\phi_4} Z_4$, we have $tpt\inv = p^2$ and $tqt\inv = q$. Since $q \in G_4$ is centralized, we have the factorization
            \[G_4 \cong \gen{p, t} \times \gen q \cong F_{20} \times Z_3.\]
            \item For $G_5 = Z_{15} \semidir_{\phi_5} Z_4$, we have $tpt\inv = p^2$ and $tqt\inv = q\inv$. Then
            \[G_5 \cong \set{p, q, t \mid p^5 = q^3 = t^4 = 1, pq = qp, tpt\inv = p^2, tqt\inv = q\inv} = F_{20} \semidir Z_3.\]
        \end{enumerate}
        \item Suppose $T = \gen u \times \gen v \cong Z_2 \times Z_2$. While this may seem like a lot of cases, observe that $\beta_i^2 = \alpha_1$ for all $i = 1, 2, 3, 4$, since either $q \mapsto q\inv \mapsto q$, or $p \mapsto p^2 \mapsto p^4 = p\inv$ ($p^3$ is handled similarly). Hence, we only need to consider the homomorphisms that send $u$ and $v$ to the the images of the $\alpha_i$'s in $\aut(Z_{15})$. Moreover, we know by \exref{5.5.6} that homomorphisms from the cyclic subgroups of $T$, i.e., $\gen u, \gen v, \gen{uv}$, to $\aut(Z_{15})$ that have conjugate (same, since $\aut(Z_{15})$ is Abelian) images give isomorphic semidirect products. For completeness, we will list all potential non-trivial homomorphisms, but group them by their images in $\aut(Z_{15})$:
        \begin{itemize}
            \item With image $\gen{\alpha_1}$, we have the three homomorphisms
            \[\theta_1 = 
            \begin{cases}
                u \mapsto \alpha_1, \\
                v \mapsto \id,
            \end{cases} \quad
            \theta_2 =
            \begin{cases}
                u \mapsto \id, \\
                v \mapsto \alpha_1,
            \end{cases} \quad
            \theta_3 =
            \begin{cases}
                u \mapsto \alpha_1, \\
                v \mapsto \alpha_1,
            \end{cases}\]
            \item With image $\gen{\alpha_2}$, we have the three homomorphisms
            \[\rho_1 = 
            \begin{cases}
                u \mapsto \alpha_2, \\
                v \mapsto \id,
            \end{cases} \quad
            \rho_2 =
            \begin{cases}
                u \mapsto \id, \\
                v \mapsto \alpha_2,
            \end{cases} \quad
            \rho_3 =
            \begin{cases}
                u \mapsto \alpha_2, \\
                v \mapsto \alpha_2,
            \end{cases}\]
            \item With image $\gen{\alpha_3}$, we have the three homomorphisms
            \[\sigma_1 = 
            \begin{cases}
                u \mapsto \alpha_3, \\
                v \mapsto \id,
            \end{cases} \quad
            \sigma_2 =
            \begin{cases}
                u \mapsto \id, \\
                v \mapsto \alpha_3,
            \end{cases} \quad
            \sigma_3 =
            \begin{cases}
                u \mapsto \alpha_3, \\
                v \mapsto \alpha_3,
            \end{cases}\]
            \item With image $V_4 = \gen{\alpha_1, \alpha_2}$ (where $\alpha_1\alpha_2 = \alpha_3$), we have the homomorphisms
            \[\tau_1 = 
            \begin{cases}
                u \mapsto \alpha_1, \\
                v \mapsto \alpha_2,
            \end{cases} \hspace{-2pt}
            \tau_2 =
            \begin{cases}
                u \mapsto \alpha_2, \\
                v \mapsto \alpha_1,
            \end{cases} \hspace{-2pt}
            \tau_3 =
            \begin{cases}
                u \mapsto \alpha_1, \\
                v \mapsto \alpha_3,
            \end{cases} \hspace{-2pt}
            \tau_4 =
            \begin{cases}
                u \mapsto \alpha_3, \\
                v \mapsto \alpha_1,
            \end{cases} \hspace{-2pt}
            \tau_5 = 
            \begin{cases}
                u \mapsto \alpha_2, \\
                v \mapsto \alpha_3,
            \end{cases} \hspace{-2pt}
            \tau_6 =
            \begin{cases}
                u \mapsto \alpha_3, \\
                v \mapsto \alpha_2,
            \end{cases}\]
        \end{itemize}
        Hence, we may consider the semidirect products corresponding to $\theta_1, \rho_1, \sigma_1$, and $\tau_1$ without loss of generality. Then we have
        \begin{enumerate}
            \item For $G_\theta = Z_{15} \semidir_{\theta_1} (Z_2 \times Z_2)$, we have $u$ acting on $Z_{15}$ by $\alpha_1$, and $v$ centralizing $Z_{15}$. Then we have the factorization $\gen v \times \gen{p, q, u}$. Observe that $q$ is centralized by $u$, and $p$ is inverted by $u$. Hence, we have
            \[G_\theta \cong \gen v \times \gen q \times \gen{p, u} \cong Z_2 \times Z_3 \times D_{10} \cong Z_6 \times D_{10}.\]
            \item For $G_\rho = Z_{15} \semidir_{\rho_1} (Z_2 \times Z_2)$, we have $u$ acting on $Z_{15}$ by $\alpha_2$, and $v$ centralizing $Z_{15}$. Then we have the factorization $\gen v \times \gen{p, q, u}$. Observe that $p$ is centralized by $u$, and $q$ is inverted by $u$. Hence, we have
            \[G_\rho \cong \gen v \times \gen p \times \gen{q, u} \cong Z_2 \times Z_5 \times D_6 \cong Z_{10} \times D_6.\]
            \item For $G_\sigma = Z_{15} \semidir_{\sigma_1} (Z_2 \times Z_2)$, we have $u$ acting on $Z_{15}$ by $\alpha_3$, and $v$ centralizing $Z_{15}$. Then we have the factorization $\gen v \times \gen{p, q, u}$. Observe that both $p$ and $q$ are inverted by $u$. Hence, we have
            \[G_\sigma \cong \gen v \times \gen{p, q, u} \cong Z_2 \times D_{30}.\]
            However, we can simplify this further by noting that $u$ inverts the element $pqv$, which has order 30. Hence, we have
            \[G_\sigma \cong \gen{pqv, u} \cong D_{60}.\]
            \item For $G_\tau = Z_{15} \semidir_{\tau_1} (Z_2 \times Z_2)$, we have $u$ acting on $Z_{15}$ by $\alpha_1$, and $v$ acting on $Z_{15}$ by $\alpha_2$. Then we have the factorization $\gen{p, q, u, v}$. Observe that $p$ is inverted by $u$ and centralized by $v$, while $q$ is inverted by $v$ and centralized by $u$. Hence, we have
            \[G_\tau \cong \gen{p, u} \times \gen{q, v} \cong D_{10} \times D_6 \cong D_{10} \times S_6. \qh\]
        \end{enumerate}
    \end{itemize}
\end{solenum}

\begin{exercise}[5.5.15]
    Let $p$ be an odd prime. Prove that every element of order 2 in $\gl_2(\f_p)$ is conjugate to a diagonal matrix with $\pm 1$'s on the diagonal. Classify the groups of order $2p^2$. [If $A$ is a $2 \times 2$ matrix with $A^2 = I$ and $v_1, v_2$ is a basis for the underlying vector space, look at $A$ acting on the vectors $w_1 = v_1 + v_2$ and $w_2 = v_1 - v_2$.]
\end{exercise}

\begin{sol}
    If $A = \pm I_2$, then $A$ is already diagonal. We now prove that if every nonzero vector of a 2-dimensional vector space is an eigenvector of $A$, then $A$ is a scalar matrix. To see this, let $e_1$ and $e_2$ be eigenvectors of $A$ corresponding to eigenvalues $\lambda_1$ and $\lambda_2$, respectively. Consider the vector $e_1 + e_2$, which is also an eigenvector of $A$ with some eigenvalue $\lambda$. Then
    \[
    A(e_1 + e_2) = \lambda (e_1 + e_2) \implies \lambda_1 e_1 + \lambda_2 e_2 = \lambda e_1 + \lambda e_2 \implies (\lambda_1 - \lambda) e_1 + (\lambda_2 - \lambda) e_2 = 0.
    \]
    Since $e_1$ and $e_2$ are linearly independent, we must have $\lambda_1 = \lambda = \lambda_2$. Hence, $A$ is a scalar matrix.

    We now take the contrapositive of this, and assume that $A$ is not a scalar matrix. Then there exists a nonzero vector $v_1$ that is not an eigenvector of $A$. Let $v_2 = Av_1$. If $v_2 = \lambda v_1$, then $Av_1 = \lambda v_1$, contradicting that $v_1$ is not an eigenvector of $A$. Hence, $v_1$ and $v_2$ are linearly independent and form a basis of the vector space.

    Let $w_1 = v_1 + v_2$ and $w_2 = v_1 - v_2$. Then
    \begin{align*}
        Aw_1 & = A(v_1 + v_2) = Av_1 + Av_2 = v_2 + A v_2 = v_2 + A^2 v_1 = v_2 + v_1 = w_1, \\
        Aw_2 & = A(v_1 - v_2) = Av_1 - Av_2 = v_2 - A v_2 = v_2 - A^2 v_1 = v_2 - v_1 = -w_2.
    \end{align*}
    Let $P = \begin{bmatrix} w_1 & w_2 \end{bmatrix}$. Then
    \[P\inv AP = 
    \begin{bmatrix}
        1 & 0 \\
        0 & -1
    \end{bmatrix}.\]
    Hence, $A$ is similar to the diagonal matrix with entries $1$ and $-1$. 

    Let $|G| = 2p^2$, $P \in \syl_p(G)$, and $Q \in \syl_2(G)$. Since $n_p \equiv 1 \bmod p$ and $n_p \mid 2$, then $p$ being odd implies that $n_p = 1$, hence $P \nsub G$. Moreover, $|P| = p^2$ implies that $P \cong Z_{p^2}$, or $Z_p \times Z_p$. Put $Q = \gen q$. We split by cases:
    \begin{itemize}
        \item The Abelian case is $Z_{p^2} \times Z_2$. If $P = \gen x \cong Z_{p^2}$, then $\aut(P) \cong Z_{p(p - 1)}$. Since $|q| = 2$ and $p$ is odd, then $p - 1$ is even and $\aut(P)$ has even order, hence it has a unique subgroup of order 2 generated by some $\alpha \in \aut(P)$, namely the inversion map. We then have $qxq\inv = x\inv$. It follows that
        \[G \cong \gen{x, q \mid x^{p^2} = q^2 = 1, qxq\inv = x\inv} \cong D_{2p^2}.\]
        \item The Abelian case is $Z_p \times Z_p \times Z_2$. If $P = \gen x \times \gen y \cong Z_p \times Z_p$, then $\aut(P) \cong \gl_2(\f_p)$. Then there are three homomorphisms from $Q$ to $\aut(P)$, namely
        \[\phi_1 = 
        \begin{cases}
            x \mapsto x, \\
            y \mapsto y
        \end{cases} \quad
        \phi_2 = 
        \begin{cases}
            x \mapsto x\inv, \\
            y \mapsto y
        \end{cases} \quad
        \phi_3 = 
        \begin{cases}
            x \mapsto x, \\
            y \mapsto y\inv
        \end{cases} \quad
        \phi_4 = 
        \begin{cases}
            x \mapsto x\inv, \\
            y \mapsto y\inv
        \end{cases} \]
        Note that $\phi_1 = \id$, hence that is the same Abelian case as before. Moreover, if we let $S$ be the exchange matrix, then $S\phi_2(q)S\inv = \phi_3(q)$. Hence the semidirect products corresponding to $\phi_2$ and $\phi_3$ are isomorphic. We then have the following non-isomorphic groups:
        \begin{itemize}
            \item $G_2 \cong (Z_p \times Z_p) \semidir_{\phi_2} Z_2$ has $y$ centralized, while $x$ is inverted. We then have $\gen y \times \gen{q, x}$, where $\gen{q, x} \cong D_{2p}$. Hence,
            \[G_2 \cong Z_p \times D_{2p}.\]
            \item $G_4 \cong (Z_p \times Z_p) \semidir_{\phi_4} Z_2$ has both $x$ and $y$ inverted by $q$. Hence, it has presentation
            \[G_4 \cong \gen{x, y, q \mid x^p = y^p = q^2 = 1, qxq\inv = x\inv, qyq\inv = y\inv} \cong (Z_p \times Z_p) \semidir Z_2. \qh\]
        \end{itemize}
    \end{itemize}
\end{sol}

\begin{exercise}[5.5.16]
    Show that there are exactly 4 distinct homomorphisms from $Z_2$ into $\aut(Z_8)$. Prove that the resulting semidirect products are the groups $Z_8 \times Z_2, D_{16}$, the quasidihedral group $QD_{16}$ (cf. \exref{2.5.11}), and the modular group $M$ (cf. \exref{2.5.14}). 
\end{exercise}

\begin{sol}
    Let $Z_8 = \gen a$ and $Z_2 = \gen b$. Note that $\aut(Z_8) \cong \units[8] = \set{1, 3, 5, 7} \cong V_4$. Hence, there are exactly 4 distinct homomorphisms from $Z_2$ into $\aut(Z_8)$, given by
    \[\phi_1 : b \mapsto a, \quad \phi_2 : b \mapsto a^3, \quad \phi_3 : b \mapsto a^5, \quad \phi_4 : b \mapsto a^7.\]
    The resulting semidirect products are:
    \begin{itemize}
        \item $G_1 \cong Z_8 \semidir_{\phi_1} Z_2 \cong Z_8 \times Z_2$, since $\phi_1$ is the trivial homomorphism.
        \item $G_2 \cong Z_8 \semidir_{\phi_2} Z_2$ has the automorphism $bab\inv = a^3$, hence $G_2 \cong QD_{16}$.
        \item $G_3 \cong Z_8 \semidir_{\phi_3} Z_2$ has the automorphism $bab\inv = a^5$, hence $G_3 \cong M$.
        \item $G_4 \cong Z_8 \semidir_{\phi_4} Z_2$ has the inversion map $bab\inv = a\inv$, hence $G_4 \cong D_{16}$.
    \end{itemize}
    By \exref{5.3.1}, we know that the three non-Abelian groups are pairwise non-isomorphic. Moreover, $Z_8 \times Z_2$ is Abelian and the others are non-Abelian, so it is also non-isomorphic to the others.
\end{sol}

\begin{exercise}[5.5.17]
    Show that for any $n \geq 3$, there are exactly 4 distinct homomorphisms from $Z_2$ into $\aut(Z_{2^n})$. Prove that the resulting semidirect groups give 4 non-isomorphic groups of order $2^{n + 1}$. [Recall \exref{2.3.21} through \exref{2.3.23}.] (These four groups together with the cyclic group and the generalized quaternion group $Q_{2^{n + 1}}$ are all the groups of order $2^{n + 1}$ which possess a cyclic subgroup of index 2.)
\end{exercise}

\begin{sol}
    Let $Z_{2^n} = \gen a$ and $Z_2 = \gen b$. Recall that $\aut(Z_{2^n}) \cong Z_2 \times Z_{2^{n - 2}}$ for $n \geq 3$. Then the elements of order dividing 2 in $\aut(Z_{2^n})$ are the elements in $\set{0, x} \times {0, y^{2^{n - 2}}}$, where $x$ is the generator of $Z_2$ and $y$ is the generator of $Z_{2^{n - 2}}$. Let $(0, 0)$ represent the identity automorphism, $(x, 0)$ represent the inversion map, $(0, y^{2^{n - 2}})$ represent the automorphism $a \mapsto a^{2^{n - 1} + 1}$, and $(x, y^{2^{n - 2}})$ represent the automorphism $a \mapsto a^{2^{n - 1} - 1}$. Then the four distinct homomorphisms from $Z_2$ into $\aut(Z_{2^n})$ are given by
    \[\phi_1 : b \mapsto (0, 0), \quad \phi_2 : b \mapsto (x, 0), \quad \phi_3 : b \mapsto (0, y^{2^{n - 2}}), \quad \phi_4 : b \mapsto (x, y^{2^{n - 2}}).\]
    Then the resulting semidirect products are:
    \begin{itemize}
        \item $G_1 \cong Z_{2^n} \semidir_{\phi_1} Z_2 \cong Z_{2^n} \times Z_2$, since $\phi_1$ is the trivial homomorphism.
        \item $G_2 \cong Z_{2^n} \semidir_{\phi_2} Z_2$ has the inversion map $bab\inv = a\inv$, hence $G_2 \cong D_{2^{n + 1}}$.
        \item $G_3 \cong Z_{2^n} \semidir_{\phi_3} Z_2$ has the automorphism $bab\inv = a^{2^{n - 1} + 1}$, hence we have the presentation
        \[G_3 \cong \gen{a, b \mid a^{2^n} = b^2 = 1, bab\inv = a^{2^{n - 1} + 1}}.\]
        \item $G_4 \cong Z_{2^n} \semidir_{\phi_4} Z_2$ has the automorphism $bab\inv = a^{2^{n - 1} - 1}$, hence we have the presentation
        \[G_4 \cong \gen{a, b \mid a^{2^n} = b^2 = 1, bab\inv = a^{2^{n - 1} - 1}}.\]
    \end{itemize}
    Since $G_1$ is Abelian, it is not isomorphic to the other three groups. To see that the other groups are not pairwise isomorphic, we begin by analyzing the centers of each group.

    We first claim that elements of the form $a^ib$ are not in the center of $G_2, G_3$, or $G_4$. To see this, observe that $a^ib$ commutes with $a$ in particular, hence $a^{i + 1}b = a^iba$ implies that $ab = ba$, or $bab\inv = a$ (where $b = b\inv$ since it has order 2). However, this only occurs in $G_1$, hence $a^ib \notin Z(G_2), Z(G_3), Z(G_4)$. We can conclude that $Z(G_2), Z(G_3), Z(G_4) \subseteq \gen a$. Note that for $a^i$ to commute with $b$, then $ba^ib\inv = a^i$, or that $a^{ri} = a^i$, where $r = -1, 2^{n - 1} + 1, 2^{n - 1} - 1$ for $G_2, G_3, G_4$, respectively. Equating exponents, it must be that $i(r - 1) \equiv 0 \bmod 2^n$. We analyze each group:
    \begin{itemize}
        \item For $G_2$, we have $r - 1 = -2$, hence $i \equiv 0 \bmod 2^{n - 1}$. Then $Z(G_2) = \gen{a^{2^{n - 1}}} \cong Z_2$. (Alternatively, we know that $D_{2^{n + 1}} = D_{2k}$ for some even $k$, and $Z(D_{2k}) = \gen{r^k} \cong Z_2$.)
        \item For $G_3$, we have $r - 1 = 2^{n - 1}$, hence $i \equiv 0 \bmod 2$. Then $Z(G_3) = \gen{a^2} \cong Z_{2^{n - 1}}$.
        \item For $G_4$, we have $r - 1 = 2^{n - 1} - 2$, hence $2i(2^{n - 2} - 1) \equiv 0 \bmod 2^n$. Since $2^{n - 2} - 1$ is odd, then $i \equiv 0 \bmod 2^{n - 1}$. Then $Z(G_4) = \gen{a^{2^{n - 1}}} \cong Z_2$.
    \end{itemize}
    We see that $Z(G_2) \cong Z(G_4)$, but $Z(G_3)$ is not isomorphic to either of them. Hence, $G_3$ is not isomorphic to $G_2$ or $G_4$. To see that $G_2$ and $G_4$ are not isomorphic, we analyze the elements of order 2 in each group. In particular, consider the elements of the form $a^ib$. Then $(a^ib)^2 = a^iba^ib = a^iba^ib\inv = a^{i + ri} = a^{i(r + 1)}$, where $r$ is as above. We analyze each group:
    \begin{itemize}
        \item For $G_2$, we have $r + 1 = 0$, hence $(a^ib)^2 = a^0 = 1$. Then all elements of the form $a^ib$ have order 2, and there are $2^n$ such elements.
        \item For $G_4$, we have $r + 1 = 2^{n - 1}$, hence $(a^ib)^2 = a^{i2^{n - 1}}$. Then $(a^ib)^2 = 1$ if and only if $i \equiv 0 \bmod 2$, hence there are $2^{n - 1}$ such elements.
    \end{itemize}
    Since $G_2$ and $G_4$ have different numbers of elements of order 2, they are not isomorphic. Hence, we have shown that the four groups are pairwise non-isomorphic.
\end{sol}

\begin{spex}[5.5.18]
    Show that if $H$ is any group, then there is a group $G$ that contains $H$ as a normal subgroup with the property that for every automorphism $\sigma$ of $H$, there is an element $g \in G$ such that conjugation by $g$ when restricted to $H$ is the given automorphism $\sigma$, i.e., every automorphism of $H$ is obtained as an inner automorphism of $G$ restricted to $H$.
\end{spex}

\begin{sol}
    Let $H$ be a group, and let $G = H \semidir \aut(H)$, where the action is given by the natural action of $\aut(H)$ on $H$. Then $H$ is a normal subgroup of $G$. Let $\sigma \in \aut(H)$. Let $g = (1, \sigma) \in G$. Then for any $(h, 1) \in \wt H$, where $\wt H \cong H$, then
    \[(1, \sigma)(h, 1)(1, \sigma)\inv = ( \sigma(h), 1).\]
    Hence conjugation by $g$ restricted to $H$ gives the automorphism $\sigma$. We remark that $G = \hol(H)$.
\end{sol}

\begin{exercise}[5.5.19]
    Let $H$ be a group of order $n$, let $K = \aut(H)$, and form $G = \hol(H) = H \semidir K$ (where $\phi$ is the identity homomorphism). Let $G$ act by left multiplication on the left cosets of $K$ in $G$, and let $\pi$ be the associated permutation representation $\pi : G \to S_n$.
    \begin{subexercises}
        \item Prove that the elements of $H$ are coset representatives for the left cosets of $K$ in $G$, and with this choice of coset representatives, $\pi$ restricted to $H$ is the regular representation of $H$.
        \item Prove $\pi(G)$ is the normalizer in $S_n$ of $\pi(H)$. Deduce that, under the regular representation of any finite group $H$ of order $n$, the normalizer in $S_n$ of the image of $H$ is isomorphic to $\hol(H)$. [Show $|G| = |N_{S_n}(\pi(H))|$ using \exref{5.5.1} and \exref{5.5.2}.]
        \item Deduce that the normalizer of the group generated by an $n$-cycle in $S_n$ is isomorphic to $\hol(Z_n)$ and has order $n\phi(n)$.
    \end{subexercises}
\end{exercise}

\begin{solenum}
    \item Recall that elements of $G$ are of the form $(h, \alpha)$, where $h \in H$ and $\alpha \in K$. Then the left cosets of $K$ in $G$ are of the form $(h, \alpha)K = (h, 1)(1, \alpha)K = (h, 1)K$. Hence, the elements of $H$ are coset representatives for the left cosets of $K$ in $G$. Let $\pi : G \to S_n$ be the associated permutation representation. Then for any $h_1, h_2 \in H$, we have
    \[\pi(h_1)(h_2K) = (h_1, 1)(h_2K) = (h_1h_2, 1)K = (h_1h_2)K.\]
    Hence, $\pi$ restricted to $H$ is the regular representation of $H$.
    \item Let $N = N_{S_n}(\pi(H))$. To see that $\pi$ is injective, let $(h, \alpha) \in \ker\pi$. Then for any $h' \in H$, we have
    \[(h, \alpha)(h'K) = (h, 1)(1, \alpha)(h'K) = (h, 1)(\alpha(h')K) = (h\alpha(h'), 1)K = h\alpha(h')K.\]
    Now let $h\alpha(h') = h'$ for all $h' \in H$. Then in particular, $h' = 1$ gives $h = 1$. Then we have $\alpha(h') = h'$ for all $h' \in H$, hence $\alpha = 1$. Therefore, $(h, \alpha) = (1, 1)$, and $\pi$ is injective.

    To see that $\pi(G) \leq N$, note that $H \nsub G$. Hence, $ghg\inv \in H$ for any $g \in G$ and $h \in H$. Consequently, $\pi(ghg\inv) \in \pi(H)$, hence $\pi(g)\pi(h)\pi(g)\inv \in \pi(H)$, and $\pi(G) \leq N$. Now let $\kappa$ be the coset $K$ in $G$. Let $N_\kappa$ denote the stabilizer of $\kappa$ in $N$. Observe, by definition, that $\pi(H) \nsub N$. From part (a), we know that $\pi(H)$ acts transitively on the left cosets of $K$ in $G$, so $N = \pi(H) \cdot N_\kappa$. Moreover, $\pi(H)$ stabilizes no coset, hence $N_\kappa \cap \pi(H) = 1$. Then $N = \pi(H) \semidir N_\kappa$, and $|N| = n|N_\kappa|$.

    Observe that conjugation by elements of $N_\kappa$ on $\pi(H)$ gives automorphisms of $\pi(H)$. Hence, we have a homomorphism $\psi : N_\kappa \to \aut(\pi(H))$. To see that $\psi$ is injective, let $g \in \ker\psi$. Then $ghg\inv = h$ for all $h \in H$, hence $g$ centralizes $\pi(H)$. Moreover, $g$ stabilizes $\kappa$, hence $g \in \pi(H)$. Then
    \[g(hK) = g(\pi(h)(\kappa)) = \pi(h)(g(\kappa)) = \pi(h)(\kappa) = hK.\]
    Hence, $g = 1$, and $\psi$ is injective. Then we have
    \[|N_\kappa| \leq |\aut(\pi(H))| = |\aut(H)|.\]
    Hence, $|N| = n|N_\kappa| \leq n|\aut(H)| = |G|$. Since $\pi$ is injective, then $|G| = |\pi(G)| \leq |N|$. Hence, $|N| = |G|$, and $\pi(G) = N$.
    \item Let $\sigma \in S_n$ be an $n$-cycle. Then $\gen\sigma \cong Z_n$ acts transitively on the set $\set{1, 2, \ldots, n}$, hence the associated permutation representation is the regular representation of $Z_n$. Since any subgroup of $S_n$ isomorphic to $Z_n$ are conjugate in $S_n$, then the normalizer of any subgroup generated by an $n$-cycle is isomorphic to the normalizer of $\gen\sigma$. By part (b), we have
    \[N_{S_n}(\gen\sigma) \cong N_{S_n}(\pi(Z_n)) \cong \hol(Z_n).\]
    Since $\aut(Z_n) \cong \units[n]$, then $|N_{S_n}(\gen\sigma)| = |\hol(Z_n)| = n\phi(n)$.
\end{solenum}

\begin{exercise}[5.5.20]
    Let $p$ be an odd prime. Prove that if $P$ is a non-cyclic $p$ group, then $P$ contains a normal subgroup $U$ with $U \cong Z_p \times Z_p$. Deduce that, for odd primes $p$, a $p$-group that contains a unique subgroup of order $p$ is cyclic. (For $p = 2$, it is a theorem that the generalized quaternion groups $Q_{2^n}$ are the only non-cyclic 2-groups which contain a unique subgroup of order 2.) [Proceed by induction on $|P|$. Let $Z$ be a subgroup of order $p$ in $Z(P)$, and let $\bar P = P/Z$. If $\bar P$ is cyclic, then $P$ is Abelian by \exref{3.1.36}---show that the result is true for Abelian groups. When $\bar P$ is not cyclic, use induction to produce a normal subgroup $\bar H$ of $\bar P$ with $\bar H \cong Z_p \times Z_p$. Let $H$ be the complete preimage of $\bar H$ in $P$, so $|H| = p^3$. Let $H_0 = \set{x \in H \mid x^p = 1}$, so that $H_0$ is a characteristic subgroup of $H$ of order $p^2$ or $p^3$ by \exref{5.4.9}. Show that a suitable subgroup of $H_0$ gives the desired normal subgroup $U$].
\end{exercise}

\begin{sol}
    The base case of $n = 2$ is clear, since $P \cong Z_p \times Z_p$ is a non-cyclic $p$-group of order $p^2$. Let $|P| = p^n$ for some $n > 2$, and assume the result is true for all non-cyclic $p$-groups of order less than $p^n$. Let $Z \leq Z(P)$ be a subgroup of order $p$, and let $\bar P = P/Z$. If $\bar P$ is cyclic, then by \exref{3.1.36}, we have that $P$ is Abelian, hence $P$ has at least 2 invariant factors. Then the subgroup generated by the elements of order $p$ in $P$ is a normal subgroup of order $p^2$, which is isomorphic to $Z_p \times Z_p$. If $\bar P$ is not cyclic, then by induction, there exists a normal subgroup $\bar H \leq \bar P$ with $\bar H \cong Z_p \times Z_p$. Let $H$ be the complete preimage of $\bar H$ in $P$, so that $|H| = p^3$. Let
    \[H_0 = \set{x \in H \mid x^p = 1}.\]
    Then $H_0$ is a characteristic subgroup of $H$ of order $p^2$ or $p^3$ by \exref{5.4.9}. If $|H_0| = p^2$, then $H_0$ is characteristic in $H$ where $H \nsub P$, so $H_0 \nsub P$, and $H_0 \cong Z_p \times Z_p$, and we can take $U = H_0$. If $|H_0| = p^3$, let $P$ act on $H$ by conjugation. Since $Z \leq Z(P)$, then we may have $P$ act on $\bar H$ by conjugation. We use the class equation, but we instead consider the action of $P$ on $\bar H$ by conjugation. Since $|\bar H| = p^2$, then the class equation gives
    \[|\bar H| = |C_{\bar H}(P)| + \sum_{i = 1}^k |P : C_P(\bar x[i])|,\]
    where $\bar x[1], \ldots, \bar x[k]$ are representatives of the non-singleton orbits of the $P$-conjugation on $\bar H$. Note that $|P : C_P(\bar x[i])|$ is a power of $p$, hence $|C_{\bar H}(P)| \equiv |\bar H| \equiv 0 \bmod p$. Since $\bar 1 \in C_{\bar H}(P)$, then $|C_{\bar H}(P)| \geq p$. Then there exists some $\bar h \in C_{\bar H}(P)$ with $\bar h \neq \bar 1$. Let $h \in H$ be the preimage of $\bar h$ in $H$. Then $h \notin Z$, and $h^p \in Z$. Let $U = \gen{h, Z}$. Then $|U| = p^2$. Since $Z \nsub P$ is clear, we observe that for any $g \in P$, then $ghg\inv \in \bar h \subseteq U$, hence $U \nsub P$. Moreover, $h^p = 1$, since $H_0 = H$, and $|h| = p$ with $h \notin Z$, hence $\gen h \cap Z = 1$. Lastly, $Z \leq Z(P)$ implies that $h$ commutes with all of $Z$. Then $U \cong Z_p \times Z_p$, and we are done.
\end{sol}

\begin{exercise}[5.5.21]
    Let $p$ be an odd prime, and let $P$ be a $p$-group. Prove that if every subgroup of $P$ is normal, then $P$ is Abelian. (Note that $Q_8$ is a non-Abelian 2-group with this property, so the result is false for $p = 2$.) [Use the preceding exercises and \exref{5.4.15}.]
\end{exercise}

\begin{sol}

\end{sol}

\begin{exercise}[5.5.22]
    Let $F$ be a field, let $n$ be a positive integer, and let $G$ be a group of upper triangular matrices in $\gl_n(F)$ (cf. \exref{2.1.16})
    \begin{subexercises}
        \item Prove that $G$ is the semidirect product $U \semidir D$, where $U$ is the set of upper triangular matrices with 1's down the diagonal (cf \exref{2.1.17}), and $D$ is the set of diagonal matrices in $\gl_n(F)$.
        \item Let $n = 2$. Recall that $U \cong F$, and $D \cong F\unt \times F\unt$ (cf \exref{3.1.11}). Describe the homomorphism from $D$ into $\aut(U)$ explicitly in terms of these isomorphisms (i.e., show how each element of $F\unt \times F\unt$ acts as an automorphism on $F$).
    \end{subexercises}
\end{exercise}

\begin{exercise}[5.5.23]
    Let $K$ and $L$ be groups, let $n$ be a positive integer, let $\rho : K \to S_n$ be a homomorphism and let $H$ be the direct product of $n$ copies of $L$. In \exref{5.1.8}, an injective homomorphism $\psi$ from $S_n$ into $\aut(H)$ was constructed by letting the elements of $S_n$ permute the $n$ factors of $H$. The composition $\psi \circ \rho$ is a homomorphism from $G$ into $\aut(H)$. The \emph{wreath product} of $L$ by $K$ is the semidirect product $H \semidir K$ with respect to this homomorphism and is denoted by $L \wr K$ (this wreath product depends on the choice of permutation representation $\rho$ of $K$---if none is given explicitly, $\rho$ is assumed to be the left regular representation of $K$).
    \begin{subexercises}
        \item Assume $K$ and $L$ are finite groups and $\rho$ is the left regular representation of $K$. Find $|L \wr K|$ in terms of $|K|$ and $|L|$.
        \item Let $p$ be a prime, let $K = L = Z_p$ and let $\rho$ be the left regular representation of $K$. Prove that $Z_p \wr Z_p$ is a non-Abelian group of order $p^{p+1}$ and is isomorphic to a Sylow $p$-subgroup of $S_{p^2}$. [The $p$ copies of $Z_p$ whose direct product makes up $H$ may be represented by $p$ disjoint $p$-cycles; these are cyclically permuted by $K$.]
    \end{subexercises}
\end{exercise}

\begin{exercise}[5.5.24]
    Let $n$ be an integer $> 1$. Prove the following classification: every group of order $n$ is Abelian if and only if $n = p_1^{\alpha_1} p_2^{\alpha_2} \cdots p_r^{\alpha_r}$, where $p_1, \ldots, p_r$ are distinct primes, $\alpha_i = 1$ or $2$ for all $i \in \set{1, \ldots, r}$ and $p_i$ does not divide $p_j^{\alpha_j} - 1$ for all $i$ and $j$. [See \exref{4.5.56}.]
\end{exercise}

\begin{exercise}[5.5.25]
    Let $H(\f_p)$ be the Heisenberg group over the finite field $\f_p = \z/p\z$ (cf. \exref{5.4.20}). Prove that $H(\f_2) \cong D_8$, and that $H(\f_p)$ has exponent $p$ and is isomorphic to the first non-Abelian group in Example 7.
\end{exercise}